\section{Companion App and Health}
\label{sec:companion-health}

\subsection{Companion App}
\label{sec:companion}

The system ships with a first-party companion app for Android and iOS that
bridges phone and desktop into a single coherent experience. It is built with
Tauri 2.0 using the same Rust backend, the same design token system, and the
same UI component library as the desktop. The companion app looks and feels
like the OS because it is built from the same parts.

The protocol is custom. \emph{KDE Connect}~\cite{kdeconnect} is a widely
used open-source protocol and application that links Linux desktops and Android
phones: clipboard sync, notification mirroring, file transfer, and remote
input. It was evaluated and rejected as the foundation for this project. Its
protocol encodes assumptions appropriate for a bolt-on integration layer,
such as per-feature capability negotiation at the session level rather than at
pairing time, that would obstruct the deeper integration features planned for
later tiers. Building on top of it would mean working around its assumptions
rather than designing for the actual requirements.

\paragraph{Feature tiers.}
Features are organized into three trust tiers established at pairing time.
Table~\ref{tab:companion-tiers} shows the breakdown. The distinction between
\texttt{paired} and \texttt{trusted} exists because Tier 2 features involve
access to graph context and health data, which require a higher trust
level than file transfer and notification mirroring. Tier 2 is restricted
to the official app; a third-party client that implements the pairing protocol
only gets Tier 1.

\begin{table}[htbp]
\caption{Companion app feature tiers}
\label{tab:companion-tiers}
\begin{tabularx}{\linewidth}{llX}
\toprule
\textbf{Tier} & \textbf{Trust level} & \textbf{Features} \\
\midrule
1 & \texttt{paired}  &
  File Drop; clipboard sync (manual trigger only, not automatic);
  notification mirroring; media playback control \\
2 & \texttt{trusted} &
  Handoff (continue reading, browsing, or editing across devices;
  project-aware when Focus Mode is active);
  Unified Timeline (cross-device activity, local only);
  graph presence (phone sends active app and content context);
  health data forwarding \\
3 & Planned          &
  Wearable data pipeline; extended context sharing \\
\bottomrule
\end{tabularx}
\end{table}

\paragraph{Protocol and transport.}
Discovery uses \emph{mDNS}~\cite{rfc6762} (Multicast DNS) service
advertisement on a shared WLAN. mDNS is the same zero-configuration discovery
protocol used by Apple's Bonjour and by Avahi on Linux; it allows devices to
announce themselves by name on a local network segment without a central DNS
server. Without a shared network, the desktop and phone locate each other via
\emph{BLE} (Bluetooth Low Energy~\cite{archwiki:bluetooth}) presence broadcast
carrying the device ID. BLE is a mode of Bluetooth designed for low-power
periodic broadcasting; a device can transmit small advertisement packets
continuously at minimal battery cost without maintaining a connection. Transport is selected by availability at connection
time, shown in Table~\ref{tab:companion-transport}.

\begin{table}[htbp]
\caption{Companion app transport selection}
\label{tab:companion-transport}
\begin{tabularx}{\linewidth}{lX}
\toprule
\textbf{Transport} & \textbf{Use case} \\
\midrule
HTTPS over LAN  & Primary transport; full feature set \\
WiFi Direct~\cite{wifidirect} & Large transfers without a shared router.
  WiFi Direct allows two devices to form a peer-to-peer 802.11 connection
  without a router or access point, at full WiFi throughput. \\
Bluetooth       & Small data only: clipboard, notifications, control commands \\
\bottomrule
\end{tabularx}
\end{table}

Pairing uses a Trust on First Use model. Both devices generate an
\emph{Ed25519}~\cite{rfc8032} keypair on first launch. Ed25519 is an elliptic-curve digital signature scheme based on Curve25519;
it produces 64-byte signatures, is fast to compute and verify, and has no
known weaknesses from weak random number generation (a historical problem with
earlier elliptic-curve schemes like ECDSA~\cite{bitcoin:ecdsa}). The first connection uses \emph{TLS}~\cite{rfc8446} (Transport Layer
Security, the protocol underlying HTTPS) with self-signed certificates rather
than certificates issued by a certificate authority. Self-signed certificates
cannot be validated by the standard CA chain, but here that is intentional: the
validation happens out-of-band via a six-digit confirmation code displayed on
the desktop and confirmed on the phone. This code is a hash of both devices'
public keys; confirming it proves that the TLS session is between the two
expected devices and not an attacker in the middle. After pairing, reconnection is automatic with no
further confirmation. Unpair deletes the key on both sides with no trace.

\paragraph{Auto-connect.}
Paired devices broadcast BLE presence in the background. When a known device
is detected, the system connects automatically using the best available
transport. No user action is required after initial pairing. Android supports
full background BLE advertising reliably. iOS imposes
\emph{CoreBluetooth Background Mode}~\cite{apple:corebluetooth} restrictions:
background BLE scanning is not permitted continuously, only in periodic
system-managed windows. Auto-connect works but is less reliable and carries a
marginally higher battery cost on the iPhone side.

\paragraph{Project-aware Handoff.}
When Focus Mode is active on the desktop (Section~\ref{sec:focus-mode}), the
Handoff feature carries project context to the phone. The companion app
receives the active project's name, description, linked resources, and recent
files from the graph, and surfaces them as a continuation card on the phone's
lock screen and in the app. Tapping the card opens the project's issue tracker
link, most recently accessed document, or notes entry depending on which was
last active. The project context is transmitted as structured metadata over
the existing Tier~2 channel; no additional pairing or permission is required
beyond the \texttt{trusted} trust level.

When Focus Mode is not active, Handoff operates in its baseline mode: the
phone receives the currently active application and document without project
context.

\paragraph{Theme synchronization.}
Because the companion app uses the same design token system as the desktop,
the active theme propagates to the phone automatically: accent color, dark and
light mode, typography. When a project accent color is active via Focus Mode,
that override propagates to the companion app as well, giving the phone the
same ambient context signal as the desktop shell. This is a structural side
effect of the shared token infrastructure, not a feature that was separately
implemented.

\paragraph{iOS limitations.}
iOS imposes constraints that cannot be worked around. There is no persistent
background BLE scanning, which makes auto-connect less reliable and requires
more manual reconnects than on Android. WiFi Direct is not available to
third-party iOS apps. USB-C wired transport is inaccessible to third-party iOS apps without
\emph{MFi} (Made for iPhone/iPad/iPod) certification~\cite{apple:mfi}, Apple's
hardware licensing program. Apps without MFi certification cannot use the
Lightning or USB-C accessory protocol for direct data transfer. Health data forwarding via \emph{HealthKit}~\cite{apple:healthkit}
(Apple's on-device health data framework, accessible to third-party apps with
explicit user consent) is feasible with explicit user permission.

Tier 1 works well on iOS. Tier 2 is functional but degraded. iOS is explicitly
a secondary platform and this limitation is communicated honestly in the app
and in documentation rather than obscured.

\subsection{Health and Wellbeing}
\label{sec:health}

The system includes a first-party health data component: a local daemon that
aggregates, stores, and optionally analyzes health and wellbeing data from
connected sources. This is not a fitness application. It is health data
infrastructure: an open, local equivalent of Apple Health or Google Fit, built
as a system component with the same privacy guarantees as the rest of the
platform.

\paragraph{Data architecture.}
Health data lives in a dedicated local store managed by a \emph{Health Daemon},
separate from the Graph Daemon. Data never leaves the device without an
explicit user action. Clinical data uses \emph{FHIR R4}~\cite{fhir:r4} (Fast Healthcare
Interoperability Resources, the international standard for structured health
data interchange maintained by HL7). FHIR represents health records as typed
REST resources such as Observation, Condition, and MedicationRequest, with
a defined JSON and XML serialization. Hospital and clinic systems across the EU
are increasingly FHIR-capable, particularly following regulatory pressure in
Germany (MIO/DiGA)~\cite{gematik:mio} and the broader eHealth interoperability
mandates~\cite{eu:ehealth}; using FHIR natively enables direct import and export
without format conversion. Non-clinical data such as steps,
sleep stages, heart rate variability, and nutrition uses an open JSON schema
with documented field definitions. There is no proprietary format; full export
is always available.

\paragraph{Data sources.}
The initial release supports manual input and phone sensors forwarded via the
Companion App: steps, GPS, and screen time. The plugin-based architecture is
designed from the start to accommodate further sources without structural
changes: wearables via the Companion App, hospital and clinic FHIR import, and
third-party applications via a documented Health API gated by the standard
permission system.

\paragraph{Wearable pipeline.}
Wearables connect to the system through the Companion App rather than directly,
using the wearable manufacturer's existing phone SDK on the phone side. The
data path is:

\begin{lstlisting}[language=bash, caption={Wearable data path}]
Wearable -> Companion App (phone, manufacturer SDK)
         -> OS Health Daemon (via Companion protocol)
\end{lstlisting}

No direct wearable protocol implementation is required in the OS. The Health
Daemon's plugin interface is designed to accept forwarded wearable data from
the start, so that adding wearable sources later does not require structural
changes. This is planned architecture, not initial release scope.

\paragraph{Graph integration.}
Health data integration with the Knowledge Graph is strictly opt-in and split
into three explicit levels, all of which default to off.
Table~\ref{tab:health-graph} lists the levels. Each level requires the
previous one; the Settings UI enforces this and does not allow incompatible
states. The AI layer's health-adjacent features are gated entirely on these
toggles. With all three off, the health module is a pure local viewer and
aggregator with no AI involvement whatsoever.

\begin{table}[htbp]
\caption{Health graph integration levels}
\label{tab:health-graph}
\begin{tabularx}{\linewidth}{lX}
\toprule
\textbf{Level} & \textbf{What it enables} \\
\midrule
1 - Health Nodes     &
  Health events appear as standalone nodes in the Knowledge Graph \\
2 - Correlations     &
  Graph links health data to activity data by time window
  (e.g.\ sleep quality correlated with focus session length;
  active focus time per project derived from Focus Mode session records) \\
3 - Pattern Recognition &
  Full cross-domain analysis across health, productivity, and
  communication data \\
\bottomrule
\end{tabularx}
\end{table}

\paragraph{Project focus time.}
Focus Mode session records (Section~\ref{sec:focus-mode}) are written to the
Knowledge Graph with project context attached. The Health Daemon reads these
records at Level~2 integration and derives a focus time metric per project per
day: the total duration of sessions where a given project was active and the
user was interacting with the system rather than idle. This metric feeds two
Health Dashboard views. The \emph{work distribution} view shows time split
across active projects over a configurable rolling window, useful for
identifying whether time allocation matches priorities. The \emph{concentration
signal} view flags extended periods of single-project focus, which correlates
with both deep work and with burnout risk depending on duration and frequency.
Neither view interprets the data for the user; both present the signal and let
the user draw conclusions.

The focus time metric requires only Level~2 graph integration; no health
sensor data is involved. It is derived entirely from session and project
activity already present in the graph.

\paragraph{Settings and permissions.}
Health has its own dedicated section in Settings rather than being buried under
general privacy settings. Controls include per-source enable and disable for
each data source independently; per-data-type enable and disable so the user
can, for example, share step data with an app but not sleep data; the graph
integration toggles from Table~\ref{tab:health-graph} presented as a grouped
sub-section with clear explanations of what each level enables; project focus
time tracking as a sub-toggle within Level~2 integration; a configurable
data retention period per data type; full FHIR and JSON export; and full
deletion with a confirmation step.

Third-party applications request health data access per data type through the
standard permission system, but health-specific confirmation dialogs are more
explicit about data sensitivity than generic permission prompts. A fitness app
requesting sleep data gets a dialog that names the data type and explains
clearly that this is health information, not a generic file access request.
